\section{第~\ref{sec:eetypes}~章　作为器件的神经元}

单个细胞的工作模式已通过电极注入电流和记录电压直接测量；积分器和分类的数学是经典理论；而大脑利用模式重构进行计算，仍是假说。

{\footnotesize
\setlength{\tabcolsep}{3pt}\renewcommand{\arraystretch}{1.15}
\begin{longtable}{>{\raggedright}p{0.5cm}>{\raggedright}p{4.0cm}>{\raggedright}p{5.6cm}>{\raggedright}p{2.3cm}>{\raggedright\arraybackslash}p{1.7cm}}
\toprule
编号 & 论断 & 实验及测量内容 & 来源 & 状态 \\
\midrule\endfirsthead
\toprule
编号 & 论断 & 实验及测量内容 & 来源 & 状态 \\
\midrule\endhead
1 & 谐振器：对抗电压变化的慢电流使神经元成为带通滤波器，峰值在几赫兹到几十赫兹（第~\ref{sec:resonator}~节） & 在脑片中向\nt{GLUT}神经元注入频率平滑上升的正弦电流并测量阻抗 $|Z(f)|$：$33^\circ$C时峰值在 $2$--$5$~Hz（$38^\circ$C时约 $7$~Hz）；谐振由慢钾电流和慢阳离子电流产生，并被持续钠电流增强。综述给出了同一手法在许多类细胞中的结果 & \cite{HuVervaekeStorm2002,HutcheonYarom2000} & 已测量 \\
2 & 慢电流的行为像电感，与膜电容一起构成振荡回路 & 测量了轴突对电流的阈下响应，并用线性化的电导方程描述；等效电路包含一个“唯象”电感，其大小依赖于电压 & \cite{Mauro1970} & 理论 \\
3 & 带反馈的群体的时间常数为 $\tau_{\text{eff}}=\tau/(1-w)$~\eqref{eq:perfectint}；$\tau=0.1$~s、$w=0.995$ 时为 $20$~s & 本章从群体的线性方程推出；作为保持眼睛位置的机制，由一篇理论工作提出 & 本章推导；\cite{Seung1996} & 理论 \\
4 & 眼位积分器靠网络而不是单个细胞的特性来保持输入 & 在鱼的眼位保持核团中进行细胞内记录：快速转动后，神经元的频率阶跃变化并保持；向单个细胞注入短暂电流只引起短暂偏移，而电压阶跃在阈下也保持——它们由突触输入维持 & \cite{Aksay2001} & 已测量 \\
5 & 无泄漏积分器需要通过学习不断调节；失调时它要么遗忘，要么饱和 & 用与眼位成 $\pm$ 比例的视觉反馈训练鱼：保持变得不稳定或出现泄漏，神经元的时间常数下降一个数量级以上，降到 $<1$~s；正常的视觉反馈逐渐恢复了稳定 & \cite{Major2004} & 已测量 \\
6 & 双稳态神经元：短暂输入开启持久平台，抑制将其关闭；这一特性由血清素开启（第~\ref{sec:bistable}~节） & 猫控制肌肉的\nt{ACET}神经元的细胞内记录：短暂的突触兴奋使神经元进入持久工作，短暂的抑制使其复位；在脊髓猫中，给予血清素前体后这种状态才出现 & \cite{Hounsgaard1988} & 已测量 \\
7 & 两类：积分器（频率从零开始增长）和谐振器（在阈值处立即出现有限频率） & 两类由分岔理论导出。实验：在皮层脑片中，\nt{GLUT}神经元能以任意小的频率开始工作（第~1型），快速\nt{GABA}神经元则从有限频率跳变开始（第~2型） & \cite{Izhikevich2007,Tateno2004} & 已测量 \\
8 & 同一个神经元可以是积分器或符合检测器：抑制缩短求和窗口 & 在海马脑片中，兴奋性输入加和到阈值的窗口短于 $2$~ms；紧随兴奋而来的前馈抑制使它缩短；在抑制较弱的树突中，窗口较宽 & \cite{Pouille2001} & 已测量 \\
9 & 由轴突构成的延迟线（表~\ref{tab:eetypes}） & 在谷仓猫头鹰中测量了通往符合检测器的轴突延迟：不同的轴突长度给出不同的延迟，检测器被调谐到声音到达的时间差 & \cite{Carr1990} & 已测量 \\
10 & 清醒时是起搏器，睡眠时是沉默的细胞 & \nt{SERO}神经元白天几乎规则地工作，在慢波睡眠中减慢，在快速眼动睡眠中几乎沉默（详见第~\ref{sec:sleep}~章） & \cite{JacobsAzmitia1992,McGintyHarper1976} & 已测量 \\
11 & 器件由离子通道以及调节它们的广播通道决定（命题~\ref{prop:eetypes}） & 已在小型网络中证明：调节物质改变神经元的固有特性和突触强度，使同一电路产生不同的节律 & \cite{Marder2012} & 已测量 \\
12 & 大脑利用模式重构进行计算（命题~\ref{prop:eetypes}） & 没有直接实验表明解决某个任务需要重构细胞的模式；只有调节物改变电路输出的实验 & \cite{Marder2012} & 假说 \\
\bottomrule
\end{longtable}}
